\section{Problem Formulation}
\label{sec:problem_formulation}

Let $q$ be a user question and let $\mathcal{E}=\{e_1,\ldots,e_N\}$ be a collection of retrievable evidence units. A base RAG system retrieves an initial evidence pool $\mathcal{E}_0 \subseteq \mathcal{E}$ and generates a draft answer $\tilde{a}$. The goal of \textsc{CoverRAG} is to transform $\tilde{a}$ into a final answer $a$ whose answer-bearing claims are covered by verifiable evidence.

\subsection{Atomic Claims and Evidence Labels}

We decompose an answer into atomic factual claims:
\begin{equation}
    \mathcal{C}(a)=\{c_1,\ldots,c_m\}.
\end{equation}
An atomic claim is a minimal factual statement that can be independently checked against evidence. For each claim $c_i$, a verifier examines a candidate evidence set $S_i \subseteq \mathcal{E}$ and returns a label
\begin{equation}
    z_i = V(c_i,S_i) \in \mathcal{Y}.
\end{equation}
In our implementation, $\mathcal{Y}$ contains four operational states:
\begin{equation}
    \mathcal{Y}=\{\mathsf{SUP},\mathsf{UNSUP},\mathsf{NOCITE},\mathsf{BADCITE}\}.
\end{equation}
$\mathsf{SUP}$ means the claim is supported by evidence. $\mathsf{UNSUP}$ means available evidence does not support it. $\mathsf{NOCITE}$ means the answer presents the claim without a usable citation. $\mathsf{BADCITE}$ means the attached citation is missing or points to evidence that does not entail the claim.

\subsection{Claim-Level Coverage State}

At iteration $t$, \textsc{CoverRAG} maintains a coverage state
\begin{equation}
    \mathcal{M}^{t}=(\mathcal{C}^{t},\mathcal{S}^{t},\mathcal{Z}^{t},\Gamma^{t},\mathbf{r}^{t}),
\end{equation}
where $\mathcal{C}^{t}$ is the current claim set, $\mathcal{S}^{t}$ stores claim-specific evidence, $\mathcal{Z}^{t}$ stores verification labels, $\Gamma^{t}$ stores citation links, and $\mathbf{r}^{t}$ stores confidence scores. The supported coverage state is
\begin{equation}
    \mathcal{C}^{+,t}=\{c_i^t:z_i^t=\mathsf{SUP}\}.
\end{equation}
This set is used as a compact representation of what the answer already covers reliably.

\subsection{Missing Answer-Unit Targets}

We define a missing answer-unit target as a concrete answer span or claim fragment that is likely useful for answering $q$ but is not covered by $\mathcal{C}^{+,t}$. Targets are generated from two sources: unsupported or citation-broken draft claims, and high-ranked evidence sentences that contain plausible answer spans not already covered by supported claims. A target is discarded if it is already covered, duplicated, generic, title-like, or irrelevant to the question.

\subsection{Objective}

The controller should improve answer utility while respecting a bounded budget of retrieval, verification, and generation calls. Conceptually, it optimizes
\begin{align}
    \max_{\pi}\quad
    &\mathrm{Correct}(a,q)
    +\lambda_1\mathrm{Recall}(a,q)
    +\lambda_2\mathrm{ClaimSupport}(a,\mathcal{E}) \nonumber\\
    &+\lambda_3\mathrm{CitationFaith}(a,\Gamma)
    -\lambda_4\mathrm{Unsupported}(a,\mathcal{E})
    -\lambda_5\mathrm{Cost}(\pi).
\end{align}
This objective is not used as a differentiable training loss. Instead, \textsc{CoverRAG} approximates it with a rule-based and verifier-guided controller that only adds content when it is both answer-useful and evidence-supported.

